% Transcription — fonds Alexandre Grothendieck, shelfmark 13, batch 2 (pages 21-40).
% Established from the facsimile archives/batches/13/batch-02.pdf (a mirror of
% https://grothendieck.umontpellier.fr/13.pdf), per the skill
% transcribe-grothendieck. The batch file carries no cover sheet and is cut from
% the folder starting at page 21: PDF page 1 is archivists' page 21, PDF page N
% is archivists' page N+20. Checked against the pencilled numbers bottom-left —
% "29", "30", "33", "34", "36", "37", "38", "39" and "40" read clearly on the
% corresponding leaves at tile resolution, and the numbers on the earlier leaves
% read at page scale; the alignment holds end to end.
%
% Pass: Opus 5 (claude-opus-5), 2026-09-19 — first pass, unchecked against the pages by a human.
%
% What the batch is. Twenty leaves, of which sixteen carry mathematics in his
% hand. Pages 21, 25, 32 and 35 are blank apart from show-through and get no
% \page{}; page 37 carries nothing but a jotting of five symbols in his hand,
% transcribed as such. The material is one subject taken up four times over, the
% same: how much of a category of motives over a field is carried by a fibre
% functor, a group, a torsor and a period. Nothing is dated and no letter falls
% in these twenty pages — the 1965 letter the inventory announces is elsewhere
% in the folder (batch 1 ends on it, pages 18-19, in English).
%
% The argument, page by page:
%
%   22-24  the abstract situation of an integrable connection on a fibre
%          functor. M a tensor category over k_0, K/k_0 an extension, a fibre
%          functor with values in k_0; the connections on V tensor K compatible
%          with tensor are the affine space over the "canonical" one (H)_0, so
%          one is a derivation of fibre functors, i.e. an element omega(x) of
%          g_K = Lie(G) tensor K. Integrability is the Maurer-Cartan identity
%          (page 23); omega lives in Omega^1_{K/k_0} tensor g, and equivalently
%          is a k_0-linear map from the dual of g to Omega^1. Page 24 closes on
%          the smallest subspace n of g whose K-points contain the omega(x):
%          u commutes with the connection as soon as it commutes with n, so one
%          wants n = g, and n may be replaced by the algebraic Lie algebra it
%          generates.
%   26-30  the torseur des périodes, the heart of the batch. k inside C, M a
%          category of motives over k; two fibre functors T_{DR} (filtered, over
%          k) and T_{B} (over Q, through Hodge structures), and the comparison
%          isomorphism (3) phi over C. Comparing T_{DR} with T_{B} tensor k
%          gives the right torsor (4) P = Isom^{tensor}(T_{B} tensor k, T_{DR})
%          under G_k, which reconstructs T_{DR} by (5), and phi becomes a point
%          phi in P(C). To carry the filtration one introduces T_{Hdg} = gr^*
%          T_{DR}, the unipotent H' inside G' = ad(P), the torsor Q of (9), then
%          H'' and G'' and the two cocharacters i_1, i_2 of (12) with their
%          relations (13). Page 27 lists the resulting data a) to g); page 28
%          reads them back and imposes conditions 1), 2), 3); page 29 adds the
%          motivic fundamental group, the element psi of Q(C) with
%          \overline{psi} = psi^{-1}, and the 2 i pi: the canonical isomorphism
%          alpha of P x^{G} G_m with G'_m x^{G'} P makes \dot{phi} equal to
%          (1/2 i pi) alpha_C, which is an arithmetic restriction on the data.
%          Page 30 is the converse: a Hodge structure with Mumford-Tate group G
%          and a k-form V_k of V_C compatible with the filtration, asked to come
%          from a motive over k; the condition is that G'_C be definable over k,
%          and for k algebraically closed that G'_C = G_C — necessary, very
%          restrictive, and certainly not sufficient (k = Q already fails).
%   31     three lines: these necessary conditions have to be completed by
%          conditions of transcendence type, and the counter-example elucidated.
%   33-34  « Situation abstraite des doubles réseaux ». M over k, two fibre
%          functors T and T' over k and K, an isomorphism phi over L; whence
%          G = Aut^{tensor}(T), the right torsor P = Isom^{tensor}(T_K, T') and
%          phi in P(L). With P trivial the datum is a class \dot{g} in
%          G(L)/G(K), the lattice T'(V) is g(V tensor K), and V inter T'(V) is
%          computed through transporters; R, the closed image of
%          (s,t) -> t s^{-1}, and the group H it generates give
%          V inter T'(V) = V^{H}, so that the functor is fully faithful exactly
%          when V^{H} = V^{G} for every V. Page 34 removes the hypothesis that
%          P be trivial, R_K being defined intrinsically by descent.
%   36     a synoptic recapitulation on one leaf: K inside C, pi_1 = Gal, A the
%          finite adeles, and the data boxed A) to G) — G over Q; the chain
%          G''_K, H''_K, Q, H'_K, G'_K, P, G_K; the absolute integrable
%          connection; u : pi_1 -> G(A); the complex picture P_C, G_C, Q_C with
%          a), b), c); a square of fundamental groups — followed by nine
%          numbered properties, from i^C_2 = i^C_1 and the factorisations of
%          pi_1(S(C)) and pi_1 through G(Q) and G(A), to the Frobenius elements
%          (« redonne Weil »), the Hodge group over R, and the Tate modules.
%   37     five symbols in his hand at the foot of an otherwise blank leaf.
%   38-40  a fresh, more careful setting-out. Page 38 introduces the
%          transcendental datum: \overline{S}_C the universal cover of S(C) at
%          xi, Gamma_G the set of homomorphisms of S = Pi_{C/R} G_{m,C} into
%          G_R above the given central i, and f : \overline{S}_C -> Gamma_G with
%          its conditions 12) a) b) c); then 13), compatibility of f(\bar{s})_C
%          with the torsor P_t and its C-trivialisation. Page 39 gives 10)
%          (granting Hodge, G^0 is generated over Q by i^C_1), the group R
%          generated by w(pi_1(S(C))) and the quotients \widetilde{G} = G/R,
%          \widetilde{P} = P/R, whose connection descends to k — statement 16) —
%          together with the descent of H', H'', G'' and the cocharacters.
%          Page 40 restarts the whole situation from a regular S of finite type
%          over Q, P = Isom^{tensor}(T_S, T_{DR}) a torsor under G_S with a
%          vertical integrable connection and a tensor filtration, and the
%          principal covering of S(C) with group G(Q); it ends on a block
%          cancelled by a long diagonal.
%
% Notation fixed for the batch, and held to. \mathbb{C}, \mathbb{Q}, \mathbb{R}
% for the fields, \mathbb{A} for the ring of finite adeles, \mathbb{G}_{m} for
% the multiplicative group and \mathbb{S} for the barred S he writes on page 38
% for Pi_{C/R} \mathbb{G}_{m,C}. \mathcal{M} is the tensor category, T, T',
% T_{B}, T_{DR}, T_{Hdg} the fibre functors, G = Aut^{tensor}(T) the group,
% G^{0} its neutral component, G' = ad(P), G'' the twist by Q, H' and H'' the
% unipotent subgroups, P and Q the torsors, R the group of pages 33-34 and 39.
% \mathfrak{g} = Lie G, \mathfrak{n} the subspace of page 24 (the reading of
% that gothic letter is flagged there and nowhere else). The caron he writes
% over Theta, omega and \mathfrak{g} on pages 22-23 is kept as \check{} — on the
% operators it marks the datum dual to the connection, on \mathfrak{g} the dual
% space. \Gamma_{G} is the space of Hodge structures on G_{\mathbb{R}}, the same
% letter batch 3 reads on pages 41-43, and \overline{S}_{\mathbb{C}},
% \widetilde{P} carry the bar and tilde he gives them. Function names — Isom,
% Aut, Hom, Transp, Rep, Gal, Im, gr — are set \operatorname; « Dér », « Modf »,
% « Str. Hodge » keep his accents in \text. Words he underlines are set \emph.
%
% What does not carry. Pages 22, 26, 27, 31, 33 and 38 are written comparatively
% fair and read almost end to end; their displayed formulas are secure. Pages
% 23, 24, 34 and 40 are quicker and lose parts of their connective prose. Pages
% 28, 29, 30, 36 and 39 are a fast palimpsest — page 30 opens under a cancelled
% boxed paragraph and is written over itself throughout, page 36 crams a whole
% recapitulation onto one leaf with three passages scribbled out past recovery,
% and page 39 cancels roughly half of what it writes. In all of them the
% displayed formulas carry and the sentences between them very largely do not.
% The apparatus is dense there in consequence, and that density is the record of
% what the leaves support, not a failure of the pass.
%
% Where the run stood at page 21. Batch 1 closes on two leaves of a letter in
% English (pages 18-19) and then on a leaf otherwise blank carrying, in his
% hand, the title of what follows: « Descriptions transcendantes de motifs en
% car. 0 et de catégories de motifs en car. 0 — doubles réseaux etc. » (page
% 20). Page 21 is blank, and the run announced by that cover sheet opens on page
% 22 with nothing carried over in the way of notation: every object is posed
% afresh. The section heading below is his, off that leaf.
%
% Batch boundary. The run does not end at page 40: the setting-out begun on page
% 38 is still in progress, page 40 breaks off inside a cancelled block, and the
% conditions a) b) c) on f are taken up and continued at the top of page 41. A
% \note at page 40 says so.

%
% Parallel pass, and what was reconciled afterwards. The folder's three batches
% were read concurrently, in three separate contexts, so batches 2 and 3 could
% not read the preceding batch's finished .tex to inherit its notation; each
% read the tail of the preceding facsimile instead. Two seams were reconciled
% in the orchestrating session once all three had landed, and nothing else was
% touched:
%   — batch 1 had set Q, R, C, Z, S, U, G_m in \mathbf and batches 2 and 3 in
%     \mathbb. Batch 1 was converted to \mathbb (205 occurrences), which is
%     both the corpus convention and the better record of the doubled strokes
%     he draws. No reading changed.
%   — the cover on page 20 had become a \section in batch 1 (closing it) and
%     again in batch 2 (opening it). It is now a heading in batch 2 alone,
%     where the run it names actually begins, and a \note in batch 1.
% Neither was a disagreement about what the leaves say.

\documentclass[11pt,a4paper]{article}
% Le préambule commun à toutes les transcriptions du fonds Grothendieck.
%
% Il tient en une idée : **l'incertitude se compose, elle ne se résout pas.**
% Une transcription qui lisse un mot douteux a détruit la seule chose pour
% laquelle on la faisait — un lecteur doit pouvoir savoir, à chaque endroit,
% ce qui était sur le feuillet et ce qui a été ajouté. Les six macros qui
% suivent sont donc l'appareil critique, et non de la décoration.
%
% Les mêmes commandes sont reconnues par `scripts/render.mjs`, qui produit la
% vue de lecture du site. Une macro ajoutée ici doit l'être aussi là-bas,
% faute de quoi le rendu échouera bruyamment — ce qui est voulu : un rendu
% silencieusement faux est pire qu'un rendu absent.

\ProvidesPackage{grothendieck}[2026/08/08 Transcriptions du fonds Grothendieck]

\usepackage{amsmath,amssymb,amsthm}
\usepackage{mathrsfs}
\usepackage{stmaryrd}
% Les diagrammes commutatifs sont partout dans ce fonds : c'est la langue
% même de la géométrie algébrique de Grothendieck, et une transcription qui
% les rendrait en prose ne transcrirait rien.
\usepackage{tikz-cd}
% Français par défaut — c'est la langue des feuillets — mais l'anglais est
% chargé aussi : la traduction et le résumé sont des documents anglais, et la
% typographie française (espaces avant les deux-points) y serait une faute.
% Ces documents basculent par \selectlanguage{english} en tête de corps.
\usepackage[english,french]{babel}
\usepackage{fontspec}
\usepackage{geometry}
\geometry{a4paper,margin=25mm,marginparwidth=28mm}
\usepackage{marginnote}
\usepackage{xcolor}
% ulem, not soul. `soul`'s \st and \ul are built on a character-by-character
% scan that XeLaTeX's font machinery breaks: the document fails with an
% undefined control sequence rather than degrading. ulem does the same two jobs
% and is Unicode-engine safe. `normalem` stops it hijacking \emph, which the
% transcriptions use for Grothendieck's own emphasis.
\usepackage[normalem]{ulem}
\usepackage{hyperref}
\hypersetup{colorlinks=true,linkcolor=black,urlcolor=black,pdfborder={0 0 0}}

% --- Métadonnées du lot -----------------------------------------------------
% Reprises telles quelles par le script de rendu, qui les affiche en tête de
% la vue de lecture. Elles fixent ce que le fichier prétend couvrir : sans
% elles, une transcription flotte sans point d'attache dans un fonds de seize
% mille feuillets.

\newcommand{\folder}[1]{\gdef\grothFolder{#1}}
\newcommand{\batch}[1]{\gdef\grothBatch{#1}}
\newcommand{\leaves}[2]{\gdef\grothFirst{#1}\gdef\grothLast{#2}}
\newcommand{\foldertitle}[1]{\gdef\grothFoldertitle{#1}}
\gdef\grothFolder{?}\gdef\grothBatch{?}\gdef\grothFirst{?}\gdef\grothLast{?}\gdef\grothFoldertitle{}

% --- Le résumé --------------------------------------------------------------
%
% Il ouvre la lecture modernisée, sous le titre « Résumé », et il est composé
% en petit. Ce n'est pas un ornement : c'est le seul passage du document qui
% ne soit pas des mathématiques, et un lecteur doit pouvoir le reconnaître
% avant de l'avoir lu — puis le sauter s'il connaît déjà le sujet, ou s'y
% arrêter s'il ne le connaît pas. Le filet qui le suit marque la reprise du
% corps.

\newenvironment{resume}
  {\small\setlength{\parskip}{0.45em}}
  {\par\medskip\hrule\medskip}

% --- Mots-clés ---------------------------------------------------------------
%
% En anglais, en clôture du résumé de la lecture modernisée : le vocabulaire
% moderne sous lequel chercher ce que les feuillets construisent. Le manifeste
% du site (`npm run manifest`) les extrait pour étiqueter la cote entière —
% cette ligne est la source unique des tags, il n'y a pas de fichier à côté.
\newcommand{\keywords}[1]{\par\smallskip\noindent
  {\footnotesize\textsc{Keywords} --- \textit{#1}}\par}

% --- L'appareil critique ----------------------------------------------------

% Le numéro de feuillet, en marge. C'est l'ancre qui permet de régler un
% désaccord sur une formule en regardant *un* feuillet plutôt qu'une chemise
% de six cents. Le site s'en sert aussi pour tourner les pages du fac-similé
% quand on fait défiler la transcription.
\newcommand{\leaf}[1]{%
  \marginnote{\footnotesize\color{gray!70}\textsf{#1}}%
  \label{leaf:#1}\ignorespaces}

% Illisible : on ne devine pas. Un mot inventé qui se lit comme les autres est
% le pire résultat possible.
\newcommand{\ill}{\textcolor{red!60!black}{[\,\ldots\,]}}

% Lecture douteuse : le mot est donné, et signalé comme incertain.
\newcommand{\uncertain}[1]{\uline{#1}}

% Ajout de l'éditeur — une lettre manquante, un mot sous-entendu.
\newcommand{\add}[1]{\textcolor{blue!50!black}{[#1]}}

% Biffé par Grothendieck lui-même. Ce qu'il a rayé fait partie du document :
% une rature dit souvent où la pensée a changé de direction.
\newcommand{\struck}[1]{\sout{#1}}

% Note du transcripteur, sur l'état matériel du feuillet.
\newcommand{\note}[1]{\par\smallskip{\small\color{gray!60!black}\textit{[#1]}}\par\smallskip}

% Note marginale de Grothendieck, distincte du corps du texte.
\newcommand{\marginal}[1]{\par\smallskip\noindent\hspace{1em}%
  {\small\color{gray!60!black}#1}\par\smallskip}

% --- Titre ------------------------------------------------------------------

\AtBeginDocument{%
  \begin{center}
    {\large\bfseries Fonds Alexandre Grothendieck --- cote n\textsuperscript{o}~\grothFolder}\\[2pt]
    {\itshape\grothFoldertitle}\\[4pt]
    {\small feuillets \grothFirst--\grothLast\ (lot \grothBatch)}\\[2pt]
    {\footnotesize\color{gray!60!black}Transcription établie d'après le fac-similé numérisé
     par l'Université de Montpellier.\\
     Les lectures incertaines sont soulignées, les illisibles notées [\,\ldots\,],
     les ajouts entre crochets.}
  \end{center}
  \vspace{1em}}

\endinput


\folder{13}
\batch{2}
\pages{21}{40}
\dating{1965}
\watermark{Édition de démonstration}
\foldertitle{Groupe de Galois motivique : notes manuscrites (s.d.), lettre (1965).}

\begin{document}

\section{Descriptions transcendantes de motifs en car.\ $0$ et de catégories de motifs en car.\ $0$ — doubles réseaux etc.}

\note{Le titre est de sa main. Il est inscrit, seul, sur un feuillet par
ailleurs blanc qui clôt le lot précédent (feuillet 20) et annonce la suite. Le
feuillet 21 est blanc ; la rédaction commence au feuillet 22, sans rien
reprendre de ce qui précède : tous les objets y sont posés à neuf.}

\page{22}

\marginal{Situation abstraite de la donnée d'une \uncertain{connexion}
\uncertain{intégrable} sur un foncteur fibre.}

\struck{\ill{}}

\note{Le haut du feuillet porte quatre lignes biffées d'un trait diagonal, où
se lisent « Dans le cas général (\ill{} pas nécessairement constants) » et
« la connaissance des opérations $\Theta_{X}$ sur les $V_{\mathbb{C}}$
($V \in \mathrm{Ob}\,\mathcal{M}$) » ; le reste ne se lit pas. La rédaction
reprend en dessous.}

Considérons une $\otimes$-catégorie $\mathcal{M}$ sur corps $k_{0}$, soit $K$
une extension de $k_{0}$. On suppose donné un foncteur fibre : valeurs dans
$k_{0}$, de $\mathcal{M}$. \uncertain{Si} ce foncteur fibre est donné par $G$
sur $k_{0}$, $\mathcal{M}$ s'identifie aux $G$-$k_{0}$-\uncertain{modules}
\ill{}. On veut déterminer les façons de mettre, sur chaque
\add{vect.\ : $K/k_{0}$} $V \otimes_{k_{0}} K$ une connexion intégrable
\struck{fonctorielle} en $V$, \add{et « compatible avec $\otimes$ »}. Or il y en
a déjà une $(H)_{0}$, déduite des structures dont on \struck{\ill{}} dispose
déjà sur chaque $V \otimes_{k_{0}} K$. Une \uncertain{telle} connexion
\struck{intégrable} sera donc de la forme

\[ \check{\Theta}_{X} = \check{\Theta}_{0X} + \check{\omega}(X) \]

\note{Le signe qu'il porte au-dessus de $\Theta$ et de $\omega$ est un accent
en forme de $\vee$ ; il est rendu ici par $\check{\ }$, et se retrouve sur
$\mathfrak{g}$ au feuillet 23.}

\[ \check{\omega} : \text{Dér}(K/k_{0}) \longrightarrow \struck{\operatorname{End}_{K}(V_{K})} \ \ K\text{-dérivations} \]

\note{« $K$-dérivations » est souligné sur le feuillet.}

\[ \check{\Theta}_{X}(\varphi) = \check{\Theta}_{0X}(\varphi) + \check{\omega}(X)(\varphi) \]

Si on veut décrire les $\check{\Theta}_{X}$ sur les $V$, de façon compatible
avec \uncertain{multiplication}, on \ill{} que $\check{\omega}(X)$, pour $V$
variable, \add{$F_{K}$} est une « dérivation » des foncteurs fibres
\struck{donnés}, donc définie par un élément de $\mathfrak{g}_{K}$, où
$\mathfrak{g} = \operatorname{Lie} G$. Soit $\omega(x) \in \mathfrak{g}_{K}$
ainsi défini.

\page{23}

on trouve donc $x \mapsto \omega(x) : \text{Dér}(K/k_{0}) \to \mathfrak{g}_{K}$,
\ill{} $K$-linéaire :

\[ \omega : \text{Dér}(K/k_{0}) \longrightarrow \mathfrak{g}_{K} \]

Si on veut avoir une connexion \emph{intégrable}, il faut écrire de plus la
condition d'intégrabilité \add{dans $(X \wedge Y)$}

\[ \Theta_{0X}\, \omega(Y) - \Theta_{0Y}\, \omega(X) \ \struck{\ill{}} \ = \ \omega([X,Y]) - [\omega(X), \omega(Y)] \]

où \struck{\ill{}} $X, Y \in \text{Dér}(K/k_{0})$ et $\Theta_{0X}$,
$\Theta_{0Y}$ \add{opèrent} sur
$\mathfrak{g}_{K} = \mathfrak{g} \otimes_{k_{0}} K$ via les \ill{} sur $K$
\ill{}.

Le plus souvent, on aura en fait

\[ \omega \in \Omega^{1}_{K/k_{0}} \otimes_{K} \mathfrak{g}_{K} \quad
\bigl( = \Omega^{1}_{K/k_{0}} \otimes_{k_{0}} \mathfrak{g} \bigr) \]

du reste que $\omega$ provient déjà d'un élément

\[ \omega \struck{_{K}} \in \Omega^{1}_{k/k_{0}} \otimes_{k} \mathfrak{g}_{k} \]

si $k$ est une ext.\ de type fini de $k_{0}$.

CAS. Si $k$ est fini, \uncertain{elle} est uniquement déterminé, si \ill{} est
en car.\ $0$). Le domaine de $\omega$, c'est \ill{} aussi : le domaine d'un
\struck{\ill{}} hom.\ $k_{0}$-linéaire

\[ \Omega = \check{\mathfrak{g}} \longrightarrow \Omega^{1}_{k/k_{0}}
\quad \bigl( \longrightarrow \Omega^{1}_{K/k_{0}} \bigr) \]

Nous allons examiner : quelles conditions, \struck{\ill{}}
\add{de $k_{0}$-\ill{}} $V, W \in \mathrm{Ob}\,\mathcal{M}$ données, \struck{les
pour} pour $u : V_{\struck{K}} \to W_{\struck{K}}$ \ill{} \struck{\ill{}} pour
que $\Theta_{x}$ ($x \in \text{Dér}(K/k_{0})$) \struck{soit} commutant : $G$,
en \struck{\ill{}} $G$ \uncertain{commence} ($k_{0}$ de car.\ $0$). Il suffit

\page{24}

d'écrire que $u$ commute : $\mathfrak{g}_{K}$, ou \ill{} hyp.\ il commute avec
$\omega(x) \in \mathfrak{g}_{K}$ ($x \in \text{Dér}(K/k_{0})$).

Soit $\mathfrak{n} \subset \mathfrak{g}$ le plus petit sous-espace vectoriel de
$\mathfrak{g}$ tel que $\mathfrak{n}_{K}$ contienne les $\omega(x)$.

\note{La lettre gothique de $\mathfrak{n}$ n'est pas assurée : le tracé
supporterait aussi bien un $\mathfrak{m}$. Elle ne reparaît pas hors de ce
feuillet.}

Alors il est clair (ou, du moins, \ill{}) que \struck{\ill{}} $u$ commute :
$\mathfrak{n}$, donc il suffit \ill{} que $\mathfrak{n} = \mathfrak{g}$ pour
pouvoir conclure (N.B.\ Il \ill{} que l'opération de la \ill{} sur
$\mathfrak{n}$ en $\mathfrak{g}$ \ill{} $\mathfrak{g}$, mais en fait on le
\struck{\ill{}} \uncertain{voit} facilement que $\mathfrak{n}$ est déjà une
sous-algèbre de Lie \ill{}). On peut \struck{\ill{}} remplacer $\mathfrak{n}$
par l'algèbre de Lie \emph{algébrique} qu'elle engendre.

\page{26}

\marginal{structure filtrée sur DR et comparaison avec Betti, Hdg \ill{}}

Soit $k \subset \mathbb{C}$, considérons une catégorie de motifs
\struck{\ill{}} \add{sur $k$}, type fini \struck{\ill{}}, $\mathcal{M}$. Il y a
deux foncteurs fibres naturels,

\[ (1) \qquad T_{DR} : \mathcal{M} \longrightarrow \text{Modf}(k) \]
\[ (2) \qquad T_{B} : \mathcal{M} \longrightarrow \struck{\ill{}}\ \text{Str.\ Hodge} \longrightarrow \text{Modf}(\mathbb{Q}) , \]

le dernier dépendant de l'\uncertain{image} $k \subset \mathbb{C}$. Le premier a
une structure de filtration. De plus, on a un isom.\ canonique

\[ (3) \qquad \varphi : T_{B} \otimes_{\mathbb{Q}} \mathbb{C} \ \overset{\sim}{\longrightarrow}\ T_{DR} \otimes_{k} \mathbb{C} \]

compatible avec filtrations, les $2$ \struck{\ill{}} de $T_{B}$.

\marginal{Soit $G$ gpe des Aut.\ de $T_{B}$.}

Comparant $T_{DR}$ à $T_{B} \otimes_{\mathbb{Q}} k$, on trouve un torseur à
droite

\[ (4) \qquad P = \operatorname{Isom}_{\otimes}(T_{B} \otimes_{\mathbb{Q}} k,\ T_{DR}) \]

sous $G_{k}$, qui permet de retrouver le foncteur $T_{DR}$ comme

\[ (5) \qquad T_{DR} = \struck{\ill{}}\ P \times^{G_{k}} (T_{B} \otimes_{\mathbb{Q}} k) . \]

De plus, \struck{\ill{}} l'isom.\ (3) définit un élément canonique

\[ \varphi \in P(\mathbb{C}) \]

\marginal{La connaissance de $G$, $P$, $\varphi$ permet de retrouver \ill{}
$\mathcal{M}$ avec (1) (2) et (3).}

Pour retrouver la structure filtrée sur $T_{DR}$, et déduire la structure de
Hodge sur $T_{B}$, i.e.\ la bigraduation sur
$T_{B} \otimes_{\mathbb{Q}} \mathbb{C}$, il convient d'introduire

\[ (6) \qquad T_{\mathrm{Hdg}} : \mathcal{M} \longrightarrow \text{Mod \struck{bi}gradué f}(k) , \qquad T_{\mathrm{Hdg}} = \operatorname{gr}^{*} T_{DR} . \]

\struck{\ill{}}

\[ (7) \qquad G'_{k} = P \times^{G_{k}} G_{k} = \operatorname{ad}(P) \simeq \operatorname{Aut}_{\otimes}(T_{DR}) \]

\struck{Aussi} \[ (8) \qquad H' \subset G'_{k} \]

est le groupe des \struck{Aut} Aut de $T_{DR}$ respectant la filtration et qui
induisent l'identité sur le gradué associé.

\page{27}

C'est un groupe unipotent, et on peut introduire

\[ (9) \qquad Q = \operatorname{Isom}_{\otimes,\ \text{ind.\ l'ident.\ sur } \operatorname{gr}^{*}}(T_{DR},\ T_{\mathrm{Hdg}}) , \]

et \ill{} que $Q$ est un torseur à droite sous $H'$. Posons

\[ (10) \qquad H'' = Q \times^{H'} H = \operatorname{ad}(H') \simeq \struck{\operatorname{Aut}}\ \text{sous-groupe de} \]

\marginal{$H'' = \operatorname{Aut}_{\otimes}(T_{\mathrm{Hdg}})$ respectant la
filtration associée à la bigraduation, et induisant l'identité sur gradué
associé \ill{}}

\[ (11) \qquad G'' = Q \times^{H'} G' = \operatorname{Aut}_{\otimes}(T_{\mathrm{Hdg}}) . \]

Mais la bigraduation de $T_{\mathrm{Hdg}}$ définit

\[ (12) \qquad i_{1}, i_{2} : \mathbb{G}_{m,k} \longrightarrow G' \]

\note{Le but de (12) est écrit $G'$ ; c'est $G''$ qu'il vient d'identifier à
$\operatorname{Aut}_{\otimes}(T_{\mathrm{Hdg}})$ en (11).}

\ill{} \struck{\ill{}}

\[ (13) \qquad \varepsilon\, i_{1}(\lambda) = \varepsilon\, i_{2}(\lambda) = \lambda , \qquad i_{1}(\lambda)\, i_{2}(\lambda) = i(\lambda) . \]

Ainsi, on obtient les données suivantes :

\begin{itemize}
  \item[a)] Gpe algébrique $G$ sur $\mathbb{Q}$, avec
  $\mathbb{G}_{m} \xrightarrow{\ i\ } G \xrightarrow{\ \varepsilon\ } \mathbb{G}_{m}$ \ill{}
  \item[b)] Groupe algébrique $G''$ sur $k$, avec
  $\mathbb{G}_{m} \xrightarrow{\ i''\ } G'' \xrightarrow{\ \varepsilon''\ } \mathbb{G}_{m}$ \ill{}
  \item[c)] hom.\ $i_{1}$, $i_{2}$ (12) satisfaisant les conditions (13)
  \item[d)] Torseur \add{à g.} $Q$ sous le groupe \add{unipotent} $H'' \subset G''$, \ill{}
\end{itemize}

On considère alors

\[ \begin{cases} H' = H'' \times^{H''} Q = \operatorname{ad}_{H''}(Q) \\
\struck{\ill{}} \quad G'' \times^{H''} Q = P \quad (\text{tors.\ à d.\ sous } G'') \\
G' = G'' \times^{G''} Q = \operatorname{ad}_{G''}(Q) \end{cases} \]

\marginal{N.B.\ $Q$ est trivial, donc $G' \simeq G''$ \ill{} groupe
canoniquement}

puis

\begin{itemize}
  \item[e)] Torseur à g.\ $P$ sous $G'$, \struck{\ill{}}
  \item[f)] Un \struck{\ill{}} Isom \ \ $G \otimes_{\mathbb{Q}} k \simeq G' \times^{G'} P = \operatorname{ad}_{G'}(P)$
  \item[g)] Un élément $\varphi \in P(\mathbb{C})$, \struck{\ill{}}
\end{itemize}

\marginal{compatible avec les structures $i$, $\varepsilon$}

\page{28}

\struck{\ill{}}

\note{Une ligne biffée ouvre le feuillet ; on y lit « Moyennant conjecture de
Hodge » et rien de plus.}

On récupère $\mathcal{M}$ comme \struck{\ill{}} $\operatorname{Rep}(G)$, $T_{B}$
(comme le foncteur \struck{\ill{}}) (cf.\ \ill{}) \add{avec structure de Hodge},
$T_{DR}$ comme \struck{$T_{B}$} $P \times^{G} (T_{B} \otimes_{\mathbb{Q}} k)$
(cf.\ e), f)), \struck{la structure de} de sorte que
$G' \simeq \operatorname{Aut}_{\otimes}(T_{DR})$. \struck{\ill{}}

Alors on a $H' \subset G'$, \ill{} $(G', G'')$ (cf.\ b), c), d)), \struck{\ill{}}
un torseur $Q$ sous $(H', H'')$, \struck{\ill{}} et
$T_{\mathrm{Hdg}} \simeq T_{DR} \times^{H'} Q$, d'où
$G'' \simeq \operatorname{Aut}_{\otimes}(T_{\mathrm{Hdg}})$ et
$i_{1}, i_{2} : \mathbb{G}_{m,k} \to G''$ définissent une bigraduation sur
$T_{\mathrm{Hdg}}$. Les \struck{valeurs propres} \ill{} définissant $H''$
\ill{} un \uncertain{tenseur} \ill{} $i_{1}$ indiquent que cette bigraduation
correspond à une \struck{\ill{}} filtration de $T_{DR}$. On en déduit une
filtration sur $T_{DR} \otimes_{k} \mathbb{C}$. Or $\varphi$ définit un isom.\
de ce dernier avec $T_{B} \otimes_{\mathbb{Q}} \mathbb{C}$, qui se trouve donc
muni d'une filtration \ill{}.

Conditions à imposer aux données a) b) c) d) e) f) g) :

\begin{enumerate}
  \item[1)] Pour t\add{ou}te repr.\ (ou une repr.\ fidèle) $V$ de $G$, la
  filtration \add{(condition \uncertain{annexe})} obtenue sur $V_{\mathbb{C}}$
  \struck{\ill{}} fait de $V$ un modèle de Hodge.
  \item[2)] Ce dernier est \emph{intrinsèque}.
  \item[3)] La restriction des scalaires de $\mathbb{C}$ à $k$ de
  $V_{\mathbb{C}}$, déduite de
  $\varphi : T_{DR} \otimes_{k} \mathbb{C} \simeq T_{B} \otimes_{\mathbb{Q}} \mathbb{C}$,
  correspond à une restriction des scalaires du motif sur $\mathbb{C}$
  définissant la structure de Hodge, qui \add{\ill{}} : b).
\end{enumerate}

N.B.\ On ne sait pas expliciter 2) et 3) \struck{\ill{}} directement en termes
des données a) \ldots\ !

On \uncertain{donne} alors, \struck{si on admet les conjectures de Hodge,}
\ill{}

\page{29}

\[ \underbrace{\pi(\operatorname{Spec} k, \xi)}_{\text{Gpe fond.\ motivique}} \longrightarrow G \]

L'image \struck{\ill{}} de la composante \ill{} doit être dans le sous-groupe
rationnel sur $\mathbb{Q}$ de $G$ engendré par
$i_{1}, i_{2} : \mathbb{G}_{m,\mathbb{C}} \to G_{\mathbb{C}}$
\struck{dérivé}.

N.B.\ La seule condition 1) permet de définir un élément

\[ \psi \in Q(\mathbb{C}) \]

par la condition que la bigraduation correspondante sur
$T_{DR} \otimes_{k} \mathbb{C} \simeq T_{B} \otimes_{\mathbb{Q}} \mathbb{C}$
satisfasse :

\[ \overline{\psi} = \psi^{-1} . \]

N.B.\ Si \add{a) à g)} proviennent d'une catégorie de motifs $\mathcal{M}$ sur
$k$, alors on a \ill{} dans

\[ P \times^{G} \mathbb{G}_{m} \ \simeq\ \mathbb{G}'_{m} \times^{G'} P \]

un isom.\ marqué $\alpha$, qui dépend \ill{} $\alpha_{\mathbb{C}}$ de
$P_{\mathbb{C}} \times^{G_{\mathbb{C}}} \mathbb{G}_{m,\mathbb{C}}$. Bien dit,
considérons $\dot{\varphi} \in (P \times^{G} \mathbb{G}_{m})(\mathbb{C})$, on
aura

\[ \dot{\varphi} = \frac{1}{2 i \pi}\, \alpha_{\mathbb{C}} \]

(à voir, que on a \ill{}), ce qui est une restriction \uncertain{arithmétique}
sur les données a) à g), sachant que $2 i \pi\, \dot{\varphi} \in$ \ill{}
rat$/k'$ \ldots

\emph{Remarquons} que les données a) à g), avec conditions 1),
\add{du type brut ; \struck{\ill{}}}, sont \struck{\ill{}} des
\add{$\otimes$-catégories} sur $\mathbb{Q}$, (munies d'un \ill{} et un
$\otimes$ foncteur \add{variables} comme \ill{}) à chaque \ill{} structure de
Hodge $(V_{\mathbb{Q}},\ V_{\mathbb{C}} = \coprod V^{pq},\ V_{k})$, i.e.\ avec
une \struck{définition} \add{variables} sur $k$ des \ill{}
$V = V_{\mathbb{Q}} \otimes_{\mathbb{Q}} \mathbb{C} = \coprod_{pq} V^{pq}$ avec
\struck{la filtration naturelle} la filtration naturelle.

\page{30}

\struck{\ill{}}

\note{Le feuillet s'ouvre sur un paragraphe encadré et biffé d'une diagonale,
où se lisent « Voyons \ill{} les données a) à g) proviennent d'une cat.\ de
motifs $\mathcal{M}$ sur $k$ » et « Admettons la conj.\ de Hodge, qui implique
que si $k$ est alg.\ clos \ldots » ; le reste ne se lit pas.}

Soit $V_{\mathbb{Q}}$, $V_{\mathbb{C}} = \sum_{pq} V^{pq}$ une structure de
Hodge, et soit $G$ le groupe de Mumford-Tate, donc
$G_{\mathbb{C}} \subset \operatorname{Aut}_{\mathbb{C}}(V)$. Soit
$V_{k} \subset V_{\mathbb{C}}$ une restriction des scalaires de $\mathbb{C}$ à
$k \subset \mathbb{C}$, \struck{\ill{}} \emph{compatible} avec la filtration
naturelle de $V_{\mathbb{C}}$. Supposons que
$(V_{\mathbb{Q}},\ V_{\mathbb{C}} = \sum V^{pq},\ V_{k} \subset V_{\mathbb{C}})$
proviennent d'un motif sur $k$. \struck{\ill{}}

\ill{} intrinsèque \ill{} ce dernier, mais : $T_{DR}$, \ill{}
\struck{conj.\ de Hodge}, un \uncertain{certain} $G'$ déf.\ $/k$,
\struck{\ill{}} doit avoir \ill{}. D'autres termes, il faut que $G'_{\mathbb{C}}$
soit définissable sur $k$ via $V_{k} \subset V_{\mathbb{C}}$ ; c'est une
condition \struck{\ill{}} ; \uncertain{faux} non triviale ! Lorsque $k$ est alg.\
clos, il faudra que $G'_{\mathbb{C}} = G_{\mathbb{C}}$, ce qui \ill{}.

Pour exprimer ceci, il faut que \add{tt} \uncertain{invariant} élément $v$ de
$\bigotimes_{\mathbb{Q}}(V' \otimes_{\mathbb{Q}} V)$ invariant sous $G$

\marginal{$\simeq \bigotimes_{k}(V'_{k} \otimes_{k} V_{k})$}

(lieu de bidegré $0$) \ill{} nécessaire dans
$\bigotimes_{k}(V'_{k} \otimes_{k} V_{k})$ ; si le premier, il \uncertain{suffit}
que $\exists$ facteurs pour $k' \subset k$, \struck{\ill{}} \ill{}
\struck{catégories} \ill{} la interprétation précédente. Ce sont \ill{} des
conditions, très restrictives sur la structure \uncertain{analogue}. Elles ne
sont certainement pas suffisantes (car pour $k = \mathbb{Q}$, $V_{k} = V_{\mathbb{Q}}$,
ce n'est en général pas possible \ill{}).

\page{31}

Il faut \uncertain{donc} compléter les conditions néc.\ précédentes par des
\struck{compatibilités} \add{conditions} du type « arg.\ de transcendance »,
bien élucider le contre-exemple précédent.

\page{33}

\marginal{Situation abstraite des doubles réseaux}

$\mathcal{M}$ $\otimes$-catégorie sur corps $k$,

$T$ foncteur fibre $/k$

$T'$ \ill{} \quad $K$ \struck{extension de $k$}

\[ \varphi : T \otimes_{k} L \ \simeq\ T' \otimes_{K} L \qquad \otimes\text{-isom.} \]

\marginal{diagramme de corps}

\begin{tikzcd}[column sep=small, row sep=small]
& L \\
K \arrow[ur, no head] \arrow[dr, no head] & \\
& k \arrow[uu, no head]
\end{tikzcd}

\note{Le diagramme des corps est tracé à la main, sans têtes de flèches ; il
est reproduit tel quel.}

Soit

\[ \begin{cases} G = \operatorname{Aut}_{\otimes}(T) & \text{gpe affine sur } k \\
P = \operatorname{Isom}_{\otimes}(T_{K}, T') & \text{torseur à droite sous } G_{K} \\
\varphi \in P(L) \end{cases} \]

Ces dernières permettant de reconstituer $\mathcal{M}$, $T$, $T'$, $\varphi$.

Supposons d'abord $P$ trivial, et soit $\xi \in P(K) \subset P(L)$. On a donc
$\varphi = \xi g$, $g \in G(L)$. Si on remplace $\xi$ par $\xi'$, on aura
$\xi' = \xi a$, $a \in G(K)$, donc $g$ remplacé par $g a$. Donc on définit
canoniquement

\[ \dot{g} \in G(L)/G(K) \]

et on voit que \ill{} maintenant les $(G, \dot{g})$ équivalent \ill{} \ill{} des
$(G, P, \varphi)$. Soit alors $V$ un $G$-module, alors le réseau
$T'(V) \subset V \otimes_{k} L$ n'est autre que $g(V \otimes_{k} K)$. On a

\[ V \cap T'(V) = \{\, x \in V \ \mid\ g^{-1} x \in V \otimes_{k} K \,\} . \]

\struck{\ill{}}

Soit $s$ l'image de $g$ dans $G_{K}$ \add{(comme sous-variétés \ill{})}
\struck{soit $y = g^{-1} x$}, $\overline{\{y\}} = T_{K} \subset G_{K}$, (alors
\struck{\ill{}}, donné par $a \in G(K)$), \ill{}
$T_{K} \subset \operatorname{Transp}_{G_{K}}(x, \gamma)$,
$y \in \operatorname{Transp}_{G_{K}}(x, y)(L)$, on a \struck{\ill{}}
\struck{$s^{-1} x = g y$} \struck{\ill{}} on a $s^{-1} a = t^{-1} x$, i.e.\
$t s^{-1} x = a$, i.e.\ $t s^{-1} \in (G_{K})_{x} = \struck{\ill{}}$.

\note{Ce bas de feuillet est écrit par-dessus lui-même et biffé ligne à ligne ;
les formules portent, les phrases qui les relient non.}

Soit donc $R$ \ill{} l'adhérence de l'image de

\[ T_{k} \times T_{k} \longrightarrow G_{K} \qquad \text{par } (s,t) \mapsto t s^{-1} ; \]

il est immédiat qu'elle \ill{} indépendant \ill{} des choix de $g$ \add{dans
$G(K)$}, \struck{\ill{}}

\[ R \subset G_{K} \]

\page{34}

l'inclusion doit avoir, si $a \in V \cap T'(W)$,

\[ R_{K} \subset G_{x,K} . \]

L'inverse est vrai, car comme $T_{K}$ est fid.\ plat sur $K$, on voit par
descente que les pts \ill{} ci-dessus \ill{} dans des $k$-alg.\
(quelconques $\Lambda$) $s \cdot x$ ($s \in T_{K}(\Lambda)$) forment \ill{} un
\ill{} bien déterminé de $V_{K}$. On sait \struck{\ill{}}

\[ S = p(R_{K}) , \qquad p : G_{K} \longrightarrow G , \]

et soit $H$ le sous-groupe algébrique \struck{fermé de} $G$ engendré par $S$.
On a donc

\[ x \in V \cap T'(V) \quad \Longleftrightarrow \quad x \in V^{H}_{\struck{\ill{}}} \]

\struck{\ill{}}

Par suite, \struck{voir que les invariants} pour tout $V$, on ait
$V \cap T'(V) = V^{G}$, i.e.\ que le foncteur
$\operatorname{Rep}(G) = \mathcal{M} \longrightarrow$ \ill{} $L(K/k)$ soit fid.\
plat, il f.\ et s.\ que pour tt $V$, \struck{les invariants} on ait

\[ V^{H} = V^{G} . \]

Supposons $P$ \struck{quelconque}. On a l'hom.\ canonique de bitorseurs

\[ P \otimes_{K} P \longrightarrow G_{K} , \qquad (s, u) \mapsto s^{-1} u \]

\struck{\ill{}}

\[ \operatorname{Spec}(L) \times_{\operatorname{Spec} K} \operatorname{Spec}(L)
\longrightarrow P \times_{\operatorname{Spec} k} P \longrightarrow G_{K} ; \]

l'image fermée $R_{K}$ est donc définie intrinsèquement. On vérifie encore (par
descente fid.\ plate) que pour \ill{}, pour le $G$-module $V$

\[ x \in V \cap T'(V) = \{\, x \in V \ \mid\ R_{K} \subset G_{x,K} \,\} . \]

Définissant $H$ comme ci-dessus, on trouve donc

\[ V \cap T'(W) = V^{H} , \]

et la \ill{} conclusion que plus haut \ldots

\note{Il écrit ici $T'(W)$ là où l'argument porte sur $T'(V)$ ; le $W$ est sur
le feuillet et se retrouve en tête de la page.}

\page{36}

$K \subset \mathbb{C}$, $\overline{K}$ clôture alg.\ dans $\mathbb{C}$,
$\pi_{1} = \pi_{1}(\operatorname{Spec} K, \overline{K}) =
\operatorname{Gal}(\overline{K}/K)$, $\mathbb{A} =$ anneau des adèles
« finis » de $\mathbb{Q}$. On \struck{considère une catégorie de motifs sur $K$},
de type fini, et le groupe de \add{$G$} Galois associé au \struck{foncteur de
Betti. Cela donne \ill{}} \add{aux structures suivantes.}

\note{Cette page est une récapitulation : il encadre sept blocs qu'il étiquette
A) à G) et les fait suivre de neuf propriétés numérotées. L'ordre suivi ici est
celui des encadrés, puis celui des numéros.}

\marginal{A) \quad $G$}

\marginal{B) \quad $G''_{K} \supset H''_{K} \cdots Q \cdots H'_{K} \subset G'_{K}
\cdots P \cdots G_{K}$, \ \ avec $i_{1}, i_{2} : \mathbb{G}_{m} \rightrightarrows G''_{K}$}

\marginal{D) \quad avec connexion absolue intégrable \ill{} invariante par
$G_{K}$ \quad [A, B structures \ill{} si $K$ abs.\ alg.]}

\marginal{C) \quad $\pi_{1} \xrightarrow{\ u\ } G(\mathbb{A})$}

\struck{$(\pi_{1} \to G/G^{0}$ ; i.e.\ $\pi_{1} \to (G/G^{0})(\mathbb{A}))$}

N.B.\ La donnée \struck{\ill{}} \ill{} implique D) ; \ill{} peut être remplacée
par la « \uncertain{densité} \ill{} » E), soumise à la condition d'être
« compatibles » avec A, B, C.

Du côté complexe :

\[ P_{\mathbb{C}} \qquad G_{\mathbb{C}} \qquad\qquad Q_{\mathbb{C}} \]

\marginal{avec pt \struck{\ill{}} $\alpha$ vect$/\mathbb{C}$, de \ill{} \ill{}
que le \ill{} \quad ; \quad avec pt $\xi_{1}$ vect$/\mathbb{C}$,}

\begin{itemize}
  \item[a)] le connexion absolue définie
  $\omega \in \Omega^{1}_{\mathbb{C}/\mathbb{Q}} \otimes_{\mathbb{Q}} \mathfrak{g}_{\mathbb{Q}}$
  \item[b)] $G'_{\struck{\ill{}}} \otimes_{K} \mathbb{C} \simeq G_{\mathbb{C}}$,
  \quad $H' \otimes_{K} \mathbb{C}$ \ \ill{} \ $H'' \otimes_{K} \mathbb{C}$ \
  \ill{} \ $G''_{\struck{1}} \otimes_{K} \mathbb{C} \simeq G_{\mathbb{C}}$
  \item[c)] d'où deux hom.\
  $i^{\mathbb{C}}_{1}, \nu^{\mathbb{C}}_{2} : \mathbb{G}_{m,\mathbb{C}} \longrightarrow G_{\mathbb{C}}$
  \item[d)] \ill{} \add{e)}, pour un modèle rigidifié convenable $S$ de $K$, de
  type fini sur \struck{\ill{}} \ill{} ($K$ de t.f.\ sur $\mathbb{Q}$)
  \[ \pi_{1}(S(\mathbb{C}), \xi) \xrightarrow{\ w\ } G(\mathbb{C}) . \]
  \add{(Le pt \ill{} identifié à l'inclusion $K \subset \mathbb{C}$)}
  \marginal{N.B.\ structure \ill{} vide si $K$ est \ill{} alg.}
  \item[e)] $P/G^{0}$ est un torseur sur $K$, de groupe $G/G^{0}$, ponctué par
  $\xi : K \hookrightarrow \mathbb{C}$, donc définit
  $v : \pi_{1} \longrightarrow (G/G^{0})(\mathbb{Q})$, i.e.\
  $\pi_{1} \to (G/G^{0})_{\mathbb{Q}}$ \ill{} on connaît $\alpha$.
\end{itemize}

\marginal{F) \quad Il en résulte aussi que le diagramme \ldots\ est commutatif.}

\begin{tikzcd}[column sep=small, row sep=small, nodes={font=\scriptsize}]
\pi_{1}(S(\mathbb{C}), \xi) \arrow[r, "u"] \arrow[d, "\mathrm{can.}"'] & G(\mathbb{Q}) \arrow[d] \\
\pi_{1}(S, \xi) \arrow[r, "v"] & G/G^{0}(\mathbb{Q})
\end{tikzcd}

\marginal{G) \quad $\pi_{1}(S(\mathbb{C}), \xi) \xrightarrow{\ w\ } G(\mathbb{C})$}

\emph{Propriétés}

\begin{enumerate}
  \item[1)] $i^{\mathbb{C}}_{2} = i^{\mathbb{C}}_{1}$ \ (donc $\alpha$ est connu
  quand on connaît \ill{})
  \item[2)] $\pi_{1}(S(\mathbb{C}), \xi) \longrightarrow G(\mathbb{C})$ se
  factorise par $G(\mathbb{Q})$,
  \item[3)] $\pi_{1} \longrightarrow G(\mathbb{A})$ se factorise par
  $\pi_{1}(S, \xi) \to G(\mathbb{A})$ pour $S$ convenable
  \item[4)] les diagr.
\end{enumerate}

\begin{tikzcd}[column sep=small, row sep=small, nodes={font=\scriptsize}]
\pi_{1}(S(\mathbb{C}), \xi) \arrow[r, "w"] \arrow[d, "\mathrm{can.}"'] & G(\mathbb{Q}) \arrow[d, "\mathrm{can.}"] \\
\pi_{1}(S, \xi) \arrow[r, "u"] & G(\mathbb{A})
\end{tikzcd}

\begin{tikzcd}[column sep=small, row sep=small, nodes={font=\scriptsize}]
\pi_{1} \arrow[r, "v"] \arrow[d, "u"'] & (G/G^{0})(\mathbb{Q}) \arrow[d] \\
G(\mathbb{A}) \arrow[r] & (G/G^{0})(\mathbb{A})
\end{tikzcd}

sont commutatifs.

\begin{enumerate}
  \item[5)] Le composé
  $\pi_{1}(S, \xi) \longrightarrow G(\mathbb{A}) \xrightarrow{\ \varepsilon\ } \mathbb{A}^{*}$
  est l'hom.\ défini par l'action de $\pi_{1}(S, \xi)$ sur les racines de
  l'unité.
  \item[6)] $v : (\pi_{1})_{\mathbb{Q}} \longrightarrow (G/G^{0})_{\mathbb{Q}}$
  un épim., i.e.\ $\pi_{1} \to (G/G^{0})(\mathbb{Q})$ un épi, et
  $G/G^{0}$ \ill{} $\mathbb{Q}$ est \emph{ct}.
  \item[7)] Pour tout point \ill{} $s$ de $S$, l'élément de Frobenius $f_{s}$ et
  \ill{} de \ill{} \ill{} dans $G(\mathbb{A})$, l'image de \ill{}
  \struck{\ill{}} de $G$ (redonne Weil \ldots).
  \item[8)] $G_{\mathbb{R}}$ muni de $i^{\mathbb{C}}_{1}$, $i^{\mathbb{C}}_{2}$
  de c) est \ill{} un groupe de Hodge.
  \item[9)] (Modules Tate) Pour tout $\ell$, l'image de $\pi_{1}$ dans
  $G(\mathbb{Q}_{\ell})$ est un sous-groupe ouvert de \struck{\ill{}}
  \ill{} \struck{\ill{}}, avec \struck{\ill{} fidèle}), et \ill{}
  est $G_{\mathbb{Q}_{\ell}}$, \ill{} $V \otimes_{\mathbb{Q}} \mathbb{Q}_{\ell}$
  \ill{} \emph{semi-simple} \ill{}.
\end{enumerate}

\marginal{devra être précisé \ldots}

\struck{\ill{}}

\page{37}

\[ G/Q \]

\[ G'_{K} \qquad P \qquad G_{K} \qquad G_{\mathbb{C}}/\mathbb{C} \]

\note{Le feuillet ne porte que ces deux lignes, de sa main, sans phrase autour.}

\page{38}

Nous \struck{allons} devons introduire une autre structure de nature
transcendante, en introduisant le rev.\ universel $\overline{S}_{\mathbb{C}}$ de
$S(\mathbb{C})$ en $\xi$ [sous-entendu, de la comp.\ connexe de $\xi$ dans
$S(\mathbb{C})$].

Soit $\Gamma_{G}$ l'ensemble des hom.\ de
$\mathbb{S} = \prod_{\mathbb{C}/\mathbb{R}} \mathbb{G}_{m,\mathbb{C}}$ dans
$G_{\mathbb{R}}$, correspondant à l'hom.\ central donné
$i : \mathbb{G}_{m} \to G$ donné par les $i(\lambda) = i_{1}(\lambda) i_{2}(\lambda)$.
On considère $\Gamma_{G}$ comme \ill{} \struck{\ill{}} la \ill{} homogène où
\struck{\ill{}} $G(\mathbb{R})$ et à fortiori $G(\mathbb{Q})$ opère. Ceci posé,
on a une \emph{application}

\[ f : \overline{S}_{\mathbb{C}} \longrightarrow \Gamma_{G} \]

satisfaisant les conditions \add{($\overline{\xi}$ pt.\ \ill{} de
$\overline{S}_{\mathbb{C}}$)}

\begin{enumerate}
  \item[12) a)] $f(\overline{\xi}) = $ l'hom.\
  $\mathbb{S} \to G_{\mathbb{R}}$ défini par $i^{\mathbb{C}}_{1}$,
  $i^{\mathbb{C}}_{2}$ de c).
  \item[b)] $f(g \overline{s}) = w(g)\, f(\overline{s})$ \quad si
  $g \in \pi_{1}(S(\mathbb{C}), \xi)$, $\overline{s} \in \overline{S}_{\mathbb{C}}$
  \item[c)] $f$ est \emph{holomorphe}.
\end{enumerate}

\begin{enumerate}
  \item[13)] (Moyennant Hodge) \struck{Pour tout $\overline{s} \in \overline{S}_{\mathbb{C}}$,}
  Le composé $\overline{S}_{\mathbb{C}} \to \Gamma_{G} \to \Gamma_{\widetilde{G}}$,
  \struck{sous-groupe algébrique de $G$} le plus grand \struck{\ill{}} qui soit
  $\mathbb{R}$-\ill{} $f(\overline{s})(\mathbb{S}) \subset G_{\mathbb{R}}$, (qui
  à priori se factorise par $S(\mathbb{C})^{\circ}$) est une application
  \emph{constante}.
\end{enumerate}

\note{Il numérote 13) deux énoncés de suite : celui qui précède et celui qui
suit. La discordance est sur le feuillet.}

\begin{enumerate}
  \item[13)] Si $\overline{s} \in \overline{S}_{\mathbb{C}}$ se trouve
  au-dessus du point $t$ de $S$ [$S$ ayant \ill{} $P$ \ill{} peut-être que $P$
  \ill{} de même \ill{} soit défini sur $S$] de \ill{} $(P_{t}, G_{k(t)})$ sur
  $k(t)$, \ill{} l'hom.
  \[ f(\overline{s})_{\mathbb{C}} : \struck{\ill{}}_{\mathbb{C}} \longrightarrow (G_{k(t)})_{\mathbb{C}} \simeq G_{\mathbb{C}} \]
  \marginal{définissant une bigraduation (donc une filtration) de \ill{} sur les
  $G$-modules}
  \ill{} $T_{\mathbb{C}}$, $T$ \ill{} obtenu sur \ill{} est \emph{compatible}
  avec le torseur $P_{t}$ sous $G_{k(t)}$ et sa $\mathbb{C}$-trivialisation
  $p(\overline{s})$ \ \add{[$p : \overline{S}_{\mathbb{C}} \to (P_{s})_{\mathbb{C}}$
  et \ill{} le $S(\mathbb{C})^{\circ}$-\ill{}, une connexion \ill{} par la
  condition que $p(\overline{\xi}) = \alpha$]}.
\end{enumerate}

\page{39}

\begin{enumerate}
  \item[10)] (Moyennant conjecture de Hodge) $G^{0}$ est engendré sur
  $\mathbb{Q}$ par $i^{\mathbb{C}}_{1}$.
\end{enumerate}

\struck{\ill{}}

\marginal{N.B.\ Ce qui précède \ill{} $K$ est \ill{} abst.\ alg.}

Soit $R \subset G$ le sous-groupe algébrique engendré par
$w(\pi_{1}(S(\mathbb{C}), \xi))$. Il est \emph{invariant} moyennant 9) [car
l'image fermée de \ill{} un sous-gpe invariant $\}$, le \ill{} la clôture alg.\
de $\pi_{1}(S(\mathbb{C}), \xi) \to \pi_{1}(S, \xi)$, \ill{} \ill{} étant
$\pi_{1}(\operatorname{Spec} K, \xi)$ \ill{} de 4), l'hom.\ \ill{} $\mathbb{Q}$
dans $K$]. \struck{\ill{}} \ill{} $\widetilde{G} = G/R$, \ill{}
$\pi_{1}(S_{\mathbb{C}}, \xi) \to G(\mathbb{A})$ \ill{}
$\widetilde{G}(\mathbb{A})$ ; \ill{} \struck{\ill{}}

\[ \widetilde{u} : \pi_{1}(\operatorname{Spec} k, \xi) \longrightarrow \widetilde{G}(\mathbb{A}) , \qquad \pi_{1}(\operatorname{Spec} k, \xi) = \pi_{1} . \]

D'autre part, on considère le torseur $\widetilde{P} = P/R$ sous
$\widetilde{G}$, avec la connexion intégrable déduite de celle de $P$. Cette
connexion intégrable « se descend à $k$ », donc (Moyennant Hodge))

\begin{enumerate}
  \item[16)] $(\widetilde{P}, \widetilde{G})$ avec sa connexion \ill{} provient
  de $(\widetilde{P}_{0}, (\widetilde{G}_{k}))$ sur $k$, et une isom.\ de
  torseurs
  \[ (*) \qquad \widetilde{P} = P/R \ \simeq\ \widetilde{P}_{0} \otimes_{k} K \]
\end{enumerate}

De plus, \struck{\ill{}} posant \struck{\ill{}}
$\widetilde{G}' = G'/R'$, de sorte que \struck{$R' \subset G'$ correspondant à
$R$}, \ill{}

\[ \widetilde{G}' = \widetilde{G}'_{0} \otimes_{k} K , \qquad
\widetilde{G}' = \operatorname{Aut}_{(\widetilde{G}_{k})}(\widetilde{P}_{0}) , \]
\[ \widetilde{H}' = H'/(H' \cap R') \ \simeq\ \operatorname{Im}(H' \to \widetilde{G}') , \]

\struck{\ill{}} « défini sur $k$ », i.e.\ le \ill{} sous-groupe \ill{} provient
de $\widetilde{H}'_{0} \subset \widetilde{G}'_{0}$. \struck{\ill{}} De plus,
pour un $\widetilde{G}$-module $M$, une structure de \ill{} \ill{} à $k$,
\ill{} c'est celle définie par descente sur $P \times^{G_{K}} M_{K}$], donnée
une \ill{} la connexion intégrable de \ill{} en \emph{filtration}, i.e.\ donnée
de descente \ill{} \emph{v.\ faible}, cette \ill{} est
\emph{horizontalement} \ill{}, \add{[donc
$\widetilde{H}' = \operatorname{Im}(H' \to \widetilde{G}')$, donc
$\widetilde{H}'_{0} \subset \widetilde{G}'$ et
$\widetilde{Q}_{1} = Q \times^{H'} \widetilde{H}'$, se descend \ill{}
$\widetilde{Q}_{0}$ ; \ill{} $\widetilde{H}'_{0}$, et (ainsi
$\widetilde{H}''_{0}$ et $\widetilde{G}''_{0}$ quand \ill{} descendent $H''$ et
$\widetilde{G}''$), et ainsi les composés
$\widetilde{i}_{1}, \widetilde{i}_{2} : \mathbb{G}_{m,K} \to \widetilde{G}'' \to \widetilde{G}''$
\ill{} descendent en
$(\widetilde{i}_{1})_{0}, (\widetilde{i}_{2})_{0} : \mathbb{G}_{m,k} \longrightarrow \widetilde{G}''_{0}$]}.

\page{40}

$S$ schéma régulier \struck{connexe} de \uncertain{t.f.} sur $\mathbb{Q}$

$k$ clôture alg.\ de $\mathbb{Q}$ dans $\Gamma(S, \mathcal{O}_{S})$, de sorte
que $k/\mathbb{Q}$ pour $S/k$ géom.\ conn.

$k \hookrightarrow \mathbb{C}$ une immersion

$\mathcal{M}$ catégorie de motifs (« iso », lisses $\tfrac{1}{2}$ simples) sur
$S$, $\mathcal{M}$ de type fini.

L'immersion $k \hookrightarrow \mathbb{C}$ définit \ill{} d'isomorphie de
foncteurs fibres \struck{rel$/\mathbb{Q}$} \ill{} $\mathcal{M}$, savoir les

\[ T_{\xi} : \mathcal{M} \longrightarrow T_{B}(M_{\xi}) , \qquad \xi \in S(\mathbb{C}) = \operatorname{Hom}_{k}(\operatorname{Spec} \mathbb{C}, S) \]

\struck{\ill{}}

Soit \struck{$G$ \ill{} un foncteur fibre \ill{} groupe alg.\ $/\mathbb{Q}$
(\ill{}) $\operatorname{Aut}_{\otimes}(T_{\xi})$} \quad
$G = \operatorname{Aut}_{\otimes}(T)$, donc
$\mathcal{M} \simeq \operatorname{Rep}(G)$.

Alors

\[ P = \operatorname{Isom}_{\otimes}(T_{S}, T_{DR}) \]

est un torseur sous $G_{S}$. Il est muni d'une \emph{connexion} intégr.\
\uncertain{verticalement} à $k$ (ou $\mathbb{Q}$, c'est kif-kif) compatible avec
$G$. De plus le foncteur $T_{DR}$, isomorphe par construction de $P$ :

\[ M \longmapsto P \times^{G_{S}} M_{S} \qquad \longrightarrow \operatorname{Rep}(G) \]

est muni d'une $\otimes$-\emph{filtration}.

\struck{Soit $\overline{S}_{\mathbb{C}}$} Considérons $S(\mathbb{C})$, et le
revêtement $\overline{P}_{\mathbb{C}}$ de $S(\mathbb{C})$ dans \ill{}
\uncertain{dessus} \ill{} $\xi \in S(\mathbb{C})$ est formé des isom.\
$T \to T_{\xi}$. Alors $\overline{S}_{\mathbb{C}} \to$ un revêtement principal
de $S(\mathbb{C})$, de groupe $G(\mathbb{Q})$, et \ill{}

\struck{\ill{}}

\note{Tout le bas du feuillet — le morphisme $p : \overline{S}_{\mathbb{C}} \to
P_{\mathbb{C}}$, la relation $p(\xi \times g) = p(\xi) \times g$, l'application
$f$ vers l'espace $\Gamma_{G}$ et les conditions a) $f$ holomorphe, b)
$f(x g) = g^{-1} f(x)$ pour $g \in G(\mathbb{Q})$, c) — est barré d'une longue
diagonale. Le texte se lit sous la rature et est donné ici comme biffé.}

\struck{$p : \overline{S}_{\mathbb{C}} \longrightarrow P_{\mathbb{C}}$}
\add{[$p(\xi \times g) = p(\xi) \times g$ ; $p$ compatible aux connexions selon
$S$-fibrés \ldots]}

\struck{$f : P_{\mathbb{C}} \to \Gamma_{G}$ \ill{} l'espace des \ill{} de Hodge
sur \ill{} $\mathbb{R}$ \ill{}, avec un $i : \mathbb{G}_{m} \to G$ fixé)}

\struck{\ill{} conditions}

\struck{Alors $f$ satisfait aux conditions}

\begin{enumerate}
  \item[a)] \struck{$f$ holomorphe}
  \item[b)] \struck{$f(x g) = g^{-1} f(x)$ \quad pour $g \in G(\mathbb{Q})$}
  \item[c)] \struck{\ill{} sur $\Gamma_{G}$ par automorphismes intérieurs}
\end{enumerate}

\marginal{N.B.\ $G(\mathbb{R})$ dans $G(\mathbb{Q})$ \ill{}}

\note{Le feuillet s'arrête ici, en cours d'énoncé : les conditions imposées à
$f$ sont reprises et poursuivies en tête du feuillet 41, au lot suivant.}

\end{document}
